\Askhsh{18557}{4}
\begin{ylh}{1}
    \item 10.2. Εμβαδόν ευθύγραμμου σχήματος - Ισοδύναμα ευθύγραμμα σχήματα
    \item 10.3. Εμβαδόν βασικών ευθύγραμμων σχημάτων.
\end{ylh}
Δίνεται τραπέζιο \text{ΑΒΓΔ} με βάσεις \text{ΑΒ} και \text{ΓΔ}, ώστε $\text{ΑΒ} > \text{ΓΔ}$. Από τις κορυφές \text{Γ} και \text{Δ} φέρουμε \text{ΓΕ}//\text{ΑΔ} και \text{ΔΖ}//\text{ΓΒ}, με \text{Ε} και \text{Ζ} σημεία στην πλευρά \text{ΑΒ} του τραπεζίου.
\begin{alist}
\item Να συγκρίνετε τα εμβαδά των τετραπλεύρων \text{ΑΔΓΕ} και \text{ΒΓΔΖ}.  \monades{9}
\item Να εκφράσετε τις περιμέτρους των τετραπλεύρων \text{ΑΔΓΕ} και \text{ΒΓΔΖ} ως συνάρτηση των πλευρών του τραπεζίου \text{ΑΒΓΔ}.  \monades{8}
\item Πώς θα πρέπει να κατασκευάσουμε το τραπέζιο \text{ΑΒΓΔ} ώστε τα τετράπλευρα \text{ΑΔΓΕ} και \text{ΒΓΔΖ} να έχουν ίσες περιμέτρους και ίσα εμβαδά;  \monades{8}
\end{alist}